\chapter{Applicazioni ed Espressività della logica}

Infine, andiamo ad analizzare come i linguaggi logici e i sistemi di deduzione possano essere applicati nell'informatica nella risoluzione di problemi. Una logica col suo calcolo di deduzione infatti fornisce un sistema formale di rappresentazione dei problemi oltre che una sistema di inferenza per risolverli. 
In tal senso, non si è molto distanti da un modello di calcolo computazionale.
Infatti, esistono linguaggi di programmazione logica, come Prolog, che si basano su questi principi per esprimere programmi e risolvere problemi.

Tra le varie applicazioni, quella su cui ci concentreremo è la specifica di problemi decisionali, ovvero problemi che possono essere risolti con una risposta sì o no.

In particolare, andremo a vedere come le sintassi e le semantiche dei linguaggi logici possano essere utilizzate per rappresentare gli input, gli output e una relazione tra di essi, in modo tale che la congiunzione delle formule risultanti sia soddisfacibile se e solo se esiste il problema da risolvere ha risposta affermativa.

In tal senso questo tipo di computazione si discosta molto dai paradigmi più noti, poiché puramente dichiarativa. Infatti, ciò che ci permette di risolvere un problema non è una serie di operazioni o manipolazioni da eseguire, ma la verifica di alcune condizioni e vincoli che legano input e output.

\section{Applicazione della logica proposizionale}

Avendo trovato un algoritmo corretto, completo e terminante per risolvere la soddisfacibilità di formule proposizionali, se è possibile riuscire a trovare una \keyword{riduzione} di un qualsiasi problema a \(\Sat\) possiamo avere un algoritmo per risolvere quel problema.

Per riduzione si intende una codifica di un'istanza qualsiasi del problema in un'istanza di \(\Sat\), ovvero una formula proposizionale, tale che la soluzione dell'istanza di \(\Sat\) ci fornisca la soluzione dell'istanza del problema originale.

Tale codifica per essere utile deve essere efficiente, ovvero deve essere eseguibile in tempo polinomiale rispetto alla dimensione dell'istanza del problema originale. Questo poiché essendo \(\Sat\) un problema NP-completo, riduzioni inefficienti porterebbero, oltre che a poter rappresentare problemi che non sono in NP, a non avere un algoritmo efficiente per risolvere il problema originale, la cui complessità è nascosta proprio nella riduzione.

Consideriamo qualche esempio semplice di riduzione a \(\Sat\), concentrandoci principalmente su problemi su grafi, che sono tra i più comuni.

\subsection{Ciclicità}
Il problema della ciclicità è il seguente: dato un grafo \(G=(V,E)\), con \(V\) insieme dei vertici e \(E\subseteq V\times V\) insieme degli archi, esiste un ciclo in \(G\)?

Un ciclo in un grafo è ovviamente definito come un percorso che contiene più di una volta lo stesso vertice, o riducendoci al caso dei cicli semplici, un percorso che inizia e finisce nello stesso vertice. Per percorso chiaramente si intende una sequenza di vertici \(v_1,v_2,\dots,v_k\) tale che \((v_i,v_{i+1})\in E\) per ogni \(i=1,\dots,k-1\).

Come già detto, per operare la riduzione, è necessario trovare una rappresentazione dell'istanza dell'input, ovvero del grafo \(G\), e dell'output, ovvero un potenziale ciclo, e dei vincoli che legano i due, in modo tale che la congiunzione di tutte le formule risultanti sia soddisfacibile se e solo se esiste un ciclo in \(G\).

Per la rappresentazione del grafo è necessario andare a rappresentare i vertici e gli archi. Essendo che \(E\subseteq V\times V\) possiamo rappresentare ogni elemento di \((u,v)\in V\times V\) con una variabile \(X_{u,v}\) e richiedere che \(X_{u,v}\) sia vera se e solo se \((u,v)\in E\).

\[\phi_E \eqdef \bigwedge_{(u,v)\in E} X_{u,v} \land \bigwedge_{(u,v)\notin E} \neg X_{u,v}\]

Tale formula intuitivamente rappresenta un assegnamento di verità alle variabili \(X_{u,v}\) che è vero se e solo se rappresenta correttamente il grafo \(G\). Di conseguenza, andando a congiungere questa formula con altre che rappresentano i vincoli che legano l'input all'output, possiamo essere certi che qualunque assegnamento soddisfi la formula risultante questo corrisponda a proprietà del grafo \(G\) in input.

Per quanto riguarda l'output, un modo alternativo alla sequenza di vertici per descrivere un ciclo è quello di rappresentarlo come sequenza di archi. Vorremo dunque mostrare che esista un \(S\subseteq E\) tale che \(S\) contiene un ciclo. Partiamo da definire le variabili per l'output come \(Y_{u,v}\) per ogni \(u,v \in V\). Dunque possiamo definire 
\[\phi_S \eqdef \bigwedge_{u,v \in V} Y_{u,v}\to X_{u,v}\]
Tale formula ci dice che ogni elemento di \(S\) deve stare in \(E\). Inoltre, per far si che \(S\) contenga un ciclo non può essere vuoto, e questo è esprimibile con
\[\phi_{\neq \emptyset} \eqdef \bigvee_{u,v\in V} Y_{u,v} \]

Ci manca solo da definire i vincoli che legano \(S\) al fatto che contenga un ciclo. Seppur in modo non estremamente intuitivo, la formula che ci permette di esprimere questo vincolo è la seguente:
\[\phi_{cyc} \eqdef \bigwedge_{u,v} (Y_{u,v \in V}\to \bigvee_{z \in V} Y_{v,z})\]

Riflettendoci, essendo che \(\phi_{\neq \emptyset}\) ci dice che esiste almeno un \(Y_{u,v}\) vero, \(\phi_{cycle}\) porta ad avere una catena di \(Y_{u,v}\) veri, che può chiudersi solo se esiste un ciclo. Infatti, l'unico modo per far si che la catena di \(Y_{u,v}\) veri non si chiuda è poter scegliere ogni volta un \(Y_{v,z}\) diverso da soddisfare, ma essendo che \(V\) è finito, prima o poi si sarà costretti a chiudere la catena, e dunque ad avere un ciclo. Equivalentemente, il percorso indotto dalle variabili di output vere per non contenere un ciclo dovrebbe essere infinito, ma essendo che \(V\) è finito, questo non è possibile. Viceversa, se esiste un ciclo, allora esiste un assegnamento che soddisfa \(\phi_{cycle}\) e \(\phi_{\not \emptyset}\) che è quello che rende vere le variabili \(Y_{u,v}\) corrispondenti agli archi del ciclo.

Dunque, concludendo, la formula che rappresenta l'istanza di \(\Sat\) a cui ridurre il problema della ciclicità è la seguente:
\[\phi \eqdef \phi_E \land \phi_S \land \phi_{\neq \emptyset} \land \phi_{cyc}\]

Per tale problema sono noti anche algoritmi polinomiali, dunque questa riduzione non è particolarmente utile, ma è un esempio semplice di come si possa rappresentare un problema con una formula proposizionale.

Passiamo ora a un altro problema, più complesso, per cui non sono noti algoritmi polinomiali, ovvero il problema della colorazione di grafi.

\subsection{K-colorazione}
Il problema della \(k\)-colorazione è il seguente: dato un grafo \(G=(V,E)\) e un intero \(k\leq \abs{V}\) numero di colori, è possibile colorare i vertici di \(G\) con al più \(k\) colori in modo tale che nessuna coppia di vertici adiacenti abbia lo stesso colore?

Più formalmente, \(\exists f:V\to C = \set{1,\dots,k} \st \forall (u,v)\in E, f(u)\neq f(v)\)?
Questo problema è molto interessante perché rappresenta una generalizzazione di un problema di allocazione di risorse. Inoltre, questo problema è molto più complicato del precedente perché si dimostra essere NP-completo.

Vediamo dunque una sua riduzione a \(\Sat\). Per rappresentare il grafo \(G\), possiamo rifarci alla rappresentazione già vista per il problema della ciclicità. Questo porta ad avere un insieme di proposizioni atomiche \(\AP[\phi_E]\) con \(\abs{\AP} = \abs{V}^2\) per rappresentare il grafo \(G\).
Per la funzione di colorazione \(f\) invece, possiamo vederla come una relazione \(f\subseteq V\times C\) con un vincolo funzionale. Dunque possiamo rappresentarla estensionalmente tramite le coppie del prodotto cartesiano.

In particolare, possiamo avere \(Y_v^c\) con \(v\in V\) e \(c\in C\), in cui ogni assegnamento di tali variabili va interpretato come una colorazione del vertice \(v\) del colore \(c\) se \(Y_v^c\) vale \(1\). Dunque \(\AP' = \set{Y_v^c}[c\in C \land v \in V]\) e quindi \(\abs{\AP'} = \abs{V}k\).

Ci rimane da imporre i vincoli per rendere tale relazione una funzione, ovvero la funzionalità e la totalità, e il vincolo che impone che vertici adiacenti non abbiano lo stesso colore.

Per la funzionalità e la totalità, ovvero che un vertice sia colorato con uno e un solo colore, possiamo imporre il seguente vincolo:
\[\phi_{f} \eqdef \bigwedge_{v\in V}\left(\bigvee_{c\in C} \left(Y_v^c \land \bigwedge_{c'\in C\setminus\set{c}} \neg Y_v^{c'}\right)\right)\]

Tale vincolo infatti impone che per ogni vertice \(v\) esista almeno un colore \(c\) tale che \(Y_v^c\) sia vero, e che per ogni colore \(c'\) diverso da \(c\), \(Y_v^{c'}\) sia falso. Dunque, ogni vertice è colorato con uno e un solo colore.

Per il vincolo che impone che vertici adiacenti non abbiano lo stesso colore, possiamo imporre il seguente vincolo:
\[\phi_{adj} \eqdef \bigwedge_{(u,v)\in E} \bigwedge_{c\in C} \neg (Y_u^c \land Y_v^c)\]

Possiamo osservare che \(\abs{\phi_{f}} = \BigO{\abs{V}k^2}\) e \(\abs{\phi_{adj}} = \BigO{\abs{V}^2k}\), dunque la formula risultante dalla congiunzione di queste due formule sarà \(\BigO{\abs{V}^3}\) essendo \(k\leq \abs{V}\). Questa codifica del problema è però problematica se l'input dei colori è rappresentato usando solo \(k\) e non l'intero insieme \(C\), poiché per rappresentare \(C\) è necessario rappresentare ogni elemento di \(C\), il che richiede spazio lineare in \(k\), mentre per rappresentare \(k\) è sufficiente rappresentare il numero \(k\) in notazione posizionale, che richiede spazio logaritmico in \(k\). In tale scenario, una lunghezza lineare in \(k\) è esponenzialmente più grande di quella dell'input e dunque la codifica non è efficiente. 

Esiste fortunatamente modo di ovviare a tale problema sfruttando proprio la rappresentazione in notazione posizionale di \(k\). In particolare, essendo che \(k\) si può rappresentare con \(\lceil \log_2(k) \rceil\) bit, è possibile rappresentare ogni colore \(c\) con una stringa di \(\lceil \log_2(k) \rceil\) bit, e dunque utilizzando \(\lceil \log_2(k) \rceil\) variabili. Dunque, invece di avere \(Y_v^c\) con \(v\in V\) e \(c\in C\), possiamo avere \(Y_v^i\) con \(v\in V\) e \(i=1,\dots,\lceil \log_2(k) \rceil\), in cui ogni assegnamento di tali variabili va interpretato come una colorazione del vertice \(v\) del colore rappresentato dalla stringa di bit formata da tali variabili. Dunque \(\AP' = \set{Y_v^i}[i=1,\dots,\lceil \log_2(k) \rceil \land v \in V]\) e quindi \(\abs{\AP'} = \BigO{\abs{V}\log k}\).

È interessante notare che questa nuova rappresentazione non necessita l'espressione del vincolo di funzionalità e totalità, poiché per ogni vertice esiste un solo insieme di variabili \(Y_v^i\) con \(i=1,\dots,\lceil \log_2(k) \rceil\) che può essere vero, e dunque ogni vertice è colorato con uno e un solo colore. Dunque, l'unico vincolo da esprimere è quello che impone che vertici adiacenti non abbiano lo stesso colore, che è il seguente:
\[\phi_2' = \bigwedge_{(u,v)\in E}\left( \bigvee_{i=0}^{\log_2(k)-1} (Y_u^i \land \neg Y_v^i) \lor (\neg Y_u^i \land Y_v^i)\right)\] 

In altre parole, quello che richiediamo adesso è che ogni coppia di vertici adiacenti abbia almeno un bit diverso nella rappresentazione del colore. Questo può essere espresso in maniera leggermente più compatta utilizzando l'operatore XOR, che è vero se e solo se i disgiunti sono diversi:
\[\phi_2' = \bigwedge_{(u,v)\in E}\left( \bigvee_{i=0}^{\log_2(k)-1} (Y_u^i \oplus Y_v^i)\right)\]

Dunque, concludendo, la formula che rappresenta l'istanza di \(\Sat\) a cui ridurre il problema della \(k\)-colorazione è la seguente:

\[\phi' \eqdef \phi_E \land \phi_2'\]

Tale codifica chiaramente occupa spazio \(\abs{V}^2\log_2(k)\), che è lineare nell'input.
Possiamo facilmente notare in oltre che in questa codifica, per \(k=2\) otteniamo tutte formule che contengono al più 2 variabili.
Tale frammento ristretto della logica proposizionale è noto come \(2-\Sat\) e per esso il problema della soddisfacibilità è risolvibile in tempo polinomiale. Dunque, questa codifica ci dice anche che \(2-\)colorazione è risolvibile in tempo polinomiale.

\subsection{Ciclo Hamiltoniano}

Un altro problema molto interessante è il problema del ciclo Hamiltoniano, che è il seguente: dato un grafo \(G=(V,E)\), esiste un ciclo semplice che contiene ogni vertice di \(G\) esattamente una volta?

Anche qui, per la rappresentazione del grafo \(G\) utilizzeremo \(\phi_E\) vista in precedenza.

Per avere un ciclo hamiltoniano, ovviamente, ci serve una sequenza di vertici lunga quanto l'insieme dei vertici, tale che ogni vertice sia presente esattamente una volta e che ci sia un arco tra ogni coppia di vertici adiacenti nella sequenza e tra l'ultimo e il primo vertice della sequenza.

Quindi per \(\abs{V}=n\) bisogna avere una sequenza di vertici \(v_0, \ldots, v_{n-1}\) tale che \(v_i \neq v_j\) per \(i\neq j\) e \((v_i,v_{i+1 \mod n})\in E\) per \(i=0,\ldots, n-1\).

È possibile vedere tale sequenza come una permutazione dell'insieme dei vertici \(V\), ovvero una funzione biettiva \(\pi:V\to \set{0,\ldots,n-1}\) che associa a ogni vertice la sua posizione nella sequenza. Dunque, per rappresentare l'output, ci basta rappresentare tale permutazione e poi imporre il vincolo che vertici adiacenti nella sequenza siano adiacenti nel grafo.

Per semplicità ci restringeremo al caso in cui \(n\) è una potenza di \(2\). Chiaramente, è possibile estendere la codifica al caso generale, ma richiede diversi accorgimenti tecnici che non sono particolarmente interessanti ai fini di questa trattazione.

Chiaramente per la rappresentazione di \(\pi\) possiamo prendere spunto da quanto fatto per la \(k\)-colorazione, e utilizzare per ogni vertice \(v\) un insieme di variabili \(Y_v^i\) con \(i=1,\dots,\log_2(n)\), in cui ogni assegnamento di tali variabili va interpretato come la rappresentazione binaria della posizione del vertice \(v\) nella sequenza. Qui l'ipotesi che \(n\) sia una potenza di \(2\) è fondamentale, poiché in questo modo è garantito che non ci possano essere vertici assegnati a posizioni che non esistono. Questo infatti garantisce che l'insieme dei vertici e l'insieme dei numeri rappresentati dalle variabili \(Y_v^i\) siano della stessa cardinalità, e dunque che sia possibile avere una corrispondenza biunivoca tra i due insiemi, che è ciò che ci serve per rappresentare una permutazione. Inoltre, questo ci permette di poter esprimere la biettività della funzione \(\pi\) utilizzando solo l'iniettività, poiché l'iniettività garantisce anche la suriettività quando i due insiemi hanno la stessa cardinalità. Dunque, per esprimere la biettività di \(\pi\) è sufficiente esprimere il vincolo di iniettività, che è il seguente:

\[\phi_{bij} \eqdef \bigwedge_{u,v\in V\st u\neq v} \bigvee_{i=0}^{\log_2(n)-1} (Y_u^i \oplus Y_v^i)\]

In altre parole, quello che tale vincolo richiede è che \(\forall u,v \in V, u\neq v \implies \pi(u)\neq\pi(v)\).
Quello che ci rimane da fare è imporre il vincolo che vertici adiacenti nella sequenza siano adiacenti nel grafo.

Iniziamo assicurandoci che tale permutazione sia effettivamente un percorso. Per fare ciò è sufficiente che ogni coppia di vertici consecutivi nella sequenza sia effettivamente un arco. Per fare ciò è sufficiente imporre che:
\[\phi_{path} \eqdef \bigwedge_{u,v\in V} \phi_{succ(u,v)} \to X_{u,v}\]
Dove \(\phi_{succ(u,v)}\) è una formula che è vera se e solo se \(\pi(v) = \pi(u)+1 \mod n\). Per definire tale formula, è necessario avere modo di codificare l'incremento di \(1\) modulo \(n\) nella rappresentazione binaria. Per fare ciò possiamo pensare a un circuito logico che prende in input la rappresentazione binaria di un numero \(x\) e restituisce in output la rappresentazione binaria di \(x+1 \mod n\). Tale circuito è composto, per ogni bit, da una porta XOR che calcola il nuovo valore del bit, e da una porta AND che calcola se è necessario portare \(1\) al bit successivo come riporto. Volendo generalizzare, è possibile vedere l'incremento di \(1\) come riporto iniziale portando a una struttura di questo tipo:

\begin{figure}[h!]
	\centering
	\begin{tikzpicture}[
			x=1cm,
			y=1cm,
			>=Latex,
			snodo/.style={draw, rectangle, minimum width=1.4cm, minimum height=0.9cm, font=\small},
			stage/.style={font=\small, align=center}
		]
		% Nodi di stadio
		\node[snodo] (b1) at (0,0) {};
		\node[snodo] (b2) at (4.2,0) {};

		% Input dall'alto
		\draw[->] (0,1.6) -- (b1.north) node[midway, right, stage] {\(x_i\)};
		\draw[->] (4.2,1.6) -- (b2.north) node[midway, right, stage] {\(x_{i+1}\)};

		% Carry in da sinistra e propagazione a destra
    \node[stage] at (-2.5,0) {\(\cdots\)};
		\draw[->] (-2.0,0) -- (b1.west) node[midway, above, stage] {\(c_i\)};
		\draw[->] (b1.east) -- (b2.west) node[midway, above, stage] {\(c_{i+1} = x_i \land c_i\)};
		\draw[->] (b2.east) -- (8,0) node[midway, above, stage] {\(c_{i+2} = x_{i+1} \land c_{i+1}\)};
		\node[stage] at (8.5,0) {\(\cdots\)};

		% Output dal basso
		\draw[->] (b1.south) -- (0,-2.0) node[midway, right, stage] {\(y_i = x_i \oplus c_i\)};
		\draw[->] (b2.south) -- (4.2,-2.0) node[midway, right, stage] {\(y_{i+1} = x_{i+1} \oplus c_{i+1}\)};

		% Carry iniziale
		\node[stage] at (-1.6,-2.4) {Incremento di 1: \(c_0 = 1\)};
	\end{tikzpicture}
	\caption{Schema a blocchi dell'incrementatore binario con propagazione del riporto.}\label{fig:binary-incrementer}
\end{figure}

Per implementare ciò sarà sufficiente introdurre le variabili per i riporto \(C_i^{u,v}\), richiedere che \(C_0^{u,v}\) sia vero, e poi imporre i vincoli che definiscono il funzionamento del circuito, ovvero \(C_{i+1}^{u,v} \leftrightarrow (Y_i^v \land C_i^{u,v})\) e \(Y_v^i \leftrightarrow (Y_u^i \oplus C_i^{u,v})\). Dunque, la formula che impone che vertici adiacenti nella sequenza siano adiacenti nel grafo è la seguente:
\[\phi_{succ(u,v)} \eqdef C_0^{u,v}\land \bigwedge_{i=0}^{\log_2(n)-1} (Y_i^u \leftrightarrow (Y_i^v \oplus C_i^{u,v}))\land (C_{i+1}^{u,v} \leftrightarrow (Y_{i}^v \land C_{i}^{u,v}))\]

A questo punto sostituendo opportunamente le varie \(\phi_{succ(u,v)}\) in \(\phi_{path}\) al variare di \(u,v\) possiamo ottenere la riduzione completa con la formula:
\[\phi = \phi_E \land \phi_{perm} \land \phi_{path}\]

\subsection{Copertura dei vertici}

Un ultimo problema che andiamo ad analizzare, anche se meno nel dettaglio, è quello della copertura dei vertici. Dato un grafo \(G=(V,E)\) e un intero \(k\leq \abs{V}\), esiste un insieme \(S\subseteq V\) di cardinalità al più \(k\) tale che ogni arco di \(E\) sia incidente ad almeno un vertice di \(S\)?

Per semplicità, analizzeremo una variante del problema che richiede il vincolo di uguaglianza sulla cardinalità di \(S\) più che la disuguaglianza. Tale semplificazione non intacca particolarmente il problema poiché se è possibile trovare una copertura di esattamente \(k\) vertici, tale copertura rispetta anche il vincolo di averne almeno \(k\).

Per la rappresentazione dell'input, come solito, possiamo rappresentare il grafo con \(\phi_E\) e \(k\) in binario.

Dobbiamo dunque definire come rappresentare l'output. Per ogni \(v \in S\) definiamo una variabile \(S_v\) la cui interpretazione sarà la relazione di appartenenza, ovvero \(S_v \iff v \in S\).

A questo punto per verificare la copertura è sufficiente imporre 

\[\phi_{cover} = \bigwedge_{(u,v) \in E} (S_v \lor S_u)\]

Tale parte è banale poiché è banale individuare una copertura qualsiasi. Infatti \(V\) è banalmente una copertura di \(G\) per definizione. Non a caso tale formula è soddisfatta assegnando a tutte le variabili \(S_v\) il valore \(1\). Il vero problema è invece quello di imporre il vincolo che \(S\) abbia esattamente \(k\) elementi.

Per fare ciò dovremmo innanzitutto usare le variabili predisposte per \(k\) per codificare una formula che sia vera se e solo se l'assegnazione di tali variabili corrisponda alla rappresentazione binaria di \(k\). Per fare ciò possiamo banalmente costruire l'insieme \(P_k\) delle posizioni della rappresentazione binaria di \(k\) che sono pari a \(1\) ed esprime il vincolo come segue:

\[\phi_k = \bigwedge_{i \in P_k} K_i \land \bigwedge_{i \notin P_k} \neg K_i\]

A questo punto per rappresentare \(\abs{S}\) abbiamo bisogno di un contatore che conti quanti elementi sono in \(S\). Per fare ciò però è sufficiente replicare la codifica dell'incrementatore visto per il ciclo hamiltoniano una volta per ogni vertice usando come riporto iniziale la variabile che rappresenta la presenza del vertice nella soluzione \(S_v\), in modo da incrementare il contatore solo se il vertice è presente. A questo punto avendo lo stato del contatore fissato a zero, ottenibile richiedendo che le variabili del primo input siano tutte false usando la negazione, confrontando l'output di ogni incrementatore con l'input del successivo in modo che siano uguali e di confrontare l'output dell'ultimo incrementatore con la rappresentazione di \(k\). Per questioni di brevità verrà omessa la rappresentazione in formule di questa costruzione, che è però facile da implementare.

\section{Espressività e complessità del primo ordine}

Passando al prim'ordine, per quanto sia possibile rappresentare problemi decisionali come nella logica proposizionale, andremo a vedere che i problemi decisionali della logica, ovvero VALIDITY e SATISFIABILITY, sono tutti indecidibili in principio.

In particolare VALIDITY è RE (ricorsivamente enumerabile) e SATISFIABILITY è coRE (il suo complemento è ricorsivamente enumerabile), ma nessuno dei due è ricorsivo. Dunque, come vedremo successivamente, esistono algoritmi che terminano in tempo finito per l'accettazione di formule valide ma possono divergere per formule non valide, e algoritmi che terminano in tempo finito per l'accettazione di formule insoddisfacibili ma possono divergere per formule soddisfacibili.

Piuttosto che approfondire quindi come ridurre problemi decisionali a VALIDITY o SATISFIABILITY, prima di mostrare come questi siano indecidibili riducendo il problema della fermata di una macchina di Turing a VALIDITY, ci concentreremo su un altro aspetto pratico interessante, ovvero la rappresentazione della conoscenza e di conseguenza l'espressività della logica del primo ordine come linguaggio.

Abbiamo già visto per la logica proposizionale alcuni problemi che potevano essere espressi usando la logica proposizionale su un input finito. In particolare tutte le proprietà finitarie erano in qualche modo esprimibili. Al prim'ordine invece possiamo esprimere proprietà con una granularità molto fine avendo a che fare con domini potenzialmente infiniti. 

Andiamo a vedere quindi alcuni esempi e tecniche molto rudimentali per dimostrare caratteristiche di espressività della logica del primo ordine. In particolare, saremo interessati a descrivere alcune proprietà di oggetti del discorso e capire se tali proprietà sono esprimibili o meno in logica del primo ordine, ovvero se è possibile trovare una formula del primo ordine che sia soddisfatta esclusivamente dai domini che soddisfano quella proprietà.

Ovviamente quest'analisi dipenderà dal linguaggio che stiamo utilizzando, ovvero dalle costanti, funzioni e predicati di cui disponiamo. 

Iniziamo definendo propriamente il concetto di proprietà e di espressibilità.

\begin{definition}{Proprietà}
  Sia \(L\) un linguaggio di \(\FOL\) e sia \(\mathcal{M}\) la classe delle strutture con \(M \in \mathcal{M}\) se \(M = \langle D,\Interp\rangle\) per \(L\). Una \textbf{proprietà} \(P\) è un qualsiasi un sott'insieme di \(\mathcal{M}\).
\end{definition}

Intuitivamente, in ottica descrittivista del linguaggio, una proprietà non è altro che l'insieme degli oggetti del discorso che la rispettano. Ad esempio, se stiamo parlando di persone, la proprietà di essere un uomo è l'insieme di tutte le persone che sono uomini.

\begin{definition}{Espressibilità}
  Sia \(L\) un linguaggio di \(\FOL\) e sia \(\mathcal{M}\) la classe delle strutture con \(M \in \mathcal{M}\) se \(M = \langle D,\Interp\rangle\) per \(L\). Una proprietà \(P\) è \textbf{esprimibile} in \(L\) se esiste una formula \(\phi \in \FOL\) tale che \(P=\den{\phi}=\set{M \in \mathcal{M} \mid M \models \phi}\)
\end{definition}

Questo problema è particolarmente interessante per l'informatica poiché ci permette di capire quanto potere espressivo sia necessario per esprimere determinate proprietà nella modellazione di basi di conoscenza. Intuitivamente l'espressività di un linguaggio è inversamente proporzionale alla complessità computazionale dei problemi di decisione di tale linguaggio, quindi è importante capire quali sono le proprietà che possiamo esprimere con un linguaggio e quali no per trovare un giusto equilibrio tra potere espressivo e complessità computazionale.

Un modo molto semplice per capire se una proprietà è esprimibile o meno è quello di cercare di costruire una formula che la esprima. 

Un esempio l'abbiamo già visto, ovvero la proprietà di essere un ordinamento lineare.
\[\phi_{lin}\eqdef\forall x, \neg (x=x) \land \forall x,y (x=y\lor x<y \lor y<x) \land \forall x,y,z ((x<y \land y<z) \to x<z)\]
con chiaramente un linguaggio che abbia come minimo un simbolo di relazione binaria \(<\) e un simbolo di uguaglianza \(=\).

Chiaramente quindi \(\den{\phi_{lin}}\) è proprio la proprietà di essere linearmente ordinato.

Sempre sullo stesso linguaggio si può pensare anche di esprimere la proprietà \(DSC = \set{M\in\mathcal{M}}[ M \text{ è ordinamento lineare discreto}]\). Sicuramente \(\phi_{lin}\) sarà parte della formula che esprime \(DSC\) ma non è sufficiente. Dovremo congiungere a \(\phi_{lin}\) anche una proprietà che esprima la discretezza. Intuitivamente quello che differenzia insiemi discreti e non discreti è il concetto di successore. Bisogna però stare attenti al fatto che un insieme discreto può ammettere massimo, e che dunque non è richiesto che ogni elemento abbia un successore. Si dovranno dunque trattare come casi particolari anche gli elementi che non hanno successore. Intuitivamente quindi quello che dovremo esprimere è che ogni elemento che ha degli elementi successivi ha un minimo elemento successivo, che è proprio il successore. La formula che esprime tutto ciò è la seguente:
\[\phi' \eqdef \forall x, ((\exists y \st x< y) \to \exists z \st(x<z \land \forall w, (x<w \to \neg(w<z))))\]
E quindi avremo che \(\den{\phi_{lin} \land \phi'} = DSC\).

Questa \(\phi_{lin} \land \phi'\) però chiaramente non caratterizza \((\N,<)\) poiché ha altri modelli che la soddisfano. Un esempio è l'unione disgiunta di \(\N\) con se stesso, ovvero \(\N + \N\), formato da coppie \((n,i)\) con \(n \in \N\) e \(i \in \set{0,1}\) e con l'ordinamento definito come segue:
\[(n,i) < (m,i) \text{ se } n < m \in \N\]
\[(n,0) < (m,1) \text{ per ogni } n,m \in \N\]

Passiamo invece alla proprietà \(DNS = \set{M\in\mathcal{M}}[ M \text{ è ordinamento lineare denso}]\). Anche in questo caso \(\phi_{lin}\) è parte della formula che esprime \(DNS\) ma non è sufficiente. Dovremo congiungere a \(\phi_{lin}\) anche
\[\phi''\eqdef \forall x,y (x<y \to \exists z (x<z \land z<y))\]

Si ha quindi \(\den{\phi_{DNS} \eqdef\phi_{lin} \land \phi'} = DNS\). Chiaramente \((\Q,<)\models \phi_{DNS}\) e \((\R,<)\models \phi_{DNS}\).

È interessante notare che \(\phi'\not\equiv\neg \phi''\). Possiamo considerare ad esempio \(([0,1]\cup \set{2},<)\). Questo chiaramente non è denso poiché per \(x = 1\) e \(y = 2\) non esiste un elemento \(z\) tale che \(1 < z < 2\). D'altra parte però non è discreto poiché per qualsiasi elemento in \([0,1]\) non esiste un successore. 

Passando a qualche risultato negativo, abbiamo che \((\Q,<)\) è indistinguibile da \((\R,<)\), ovvero non esiste una formula \(\phi\) tale che \((\Q,<)\models \phi\) e \((\R,<)\not\models \phi\). Questa è conseguenza diretta del teorema di Löwenheim-Skolem che recita quanto segue:
\begin{theorem}{Löwenheim-Skolem}
  Se \(\Gamma\) è un insieme di formule chiuse al più numerabile e soddisfacibile, allora \(\Gamma\) ha un modello al più numerabile.
\end{theorem}

Sia infatti \(\phi\) la potenziale formula che distingue \((\Q,<)\) da \((\R,<)\). Dunque \((\Q,<)\models \phi\) e \((\R,<)\not\models \phi\). 
Equivalentemente, \((\R,<)\models \neg \phi\) e \((\Q,<)\not\models \neg \phi\). Dunque \(\set{\neg \phi, \phi_{DNS}}\) è un insieme di formule chiuse al più numerabile e soddisfatto da \((\R,<)\). Dal teorema di Löwenhaim-Skolem esiste però un modello al più numerabile che soddisfa \(\set{\neg \phi, \phi_{DNS}}\). Tale modello non può essere \((\R,<)\) poiché non è al più numerabile, quindi deve essere \((\Q,<)\). Dunque \((\Q,<)\models \neg \phi\) e \((\Q,<)\not\models \neg \phi\), il che è assurdo. 

Intuitivamente il teorema di Löwenhaim-Skolem vale per la costruzione vista per il teorema di completezza, poiché la costruzione del modello dell'insieme soddisfacibile \(\Gamma\) portava ad avere un dominio al più numerabile, poiché abbiamo aggiunto un numero enumerabile di insiemi enumerabile di costanti che è ancora un insieme enumerabile se lo era quello di partenza.

Questo quindi ci dice che ci sono limiti in ciò che il prim'ordine può esprimere.

\begin{theorem}{Compattezza}
  Per ogni \(\Gamma\subseteq\FOL\) (formule chiuse) e per ogni formula \(\phi\) (chiusa), vale che
  
  \[\Gamma \models \phi \iff \text{ esiste } \Gamma_0 \subseteq \Gamma \text{ finito tale che } \Gamma_0 \models \phi\]
\end{theorem}

\begin{proof}
  Si ha per completezza che \(\Gamma \models \phi \implies \Gamma \vdash \phi\). Se \(\Gamma \vdash \phi\) allora esiste una dimostrazione di \(\phi\) a partire da \(\Gamma\) che è finita, quindi coinvolge solo un numero finito di formule di \(\Gamma\), ovvero esiste un sottoinsieme finito \(\Gamma_0\) di \(\Gamma\) tale che \(\Gamma_0 \vdash \phi\). Per correttezza, se \(\Gamma_0 \vdash \phi\) allora \(\Gamma_0 \models \phi\).

  L'altra direzione è banale per monotònicità
\end{proof}

Una formulazione alternativa della compattezza è la seguente:
\begin{theorem}{Compattezza Alternativa}
  Per ogni \(\Gamma\subseteq\FOL\) (formule chiuse),

  \[\Gamma \in\Unsat \iff \text{ esiste } \Gamma_0 \subseteq \Gamma \text{ finito tale che } \Gamma_0 \in\Unsat\]
  o equivalentemente se
  \[\Gamma \in\Sat \iff \text{ per ogni } \Gamma_0 \subseteq \Gamma \text{ finito, } \Gamma_0 \in\Sat\]
\end{theorem}

Il teorema di compattezza è uno strumento particolarmente usato, soprattutto in logica matematica, per dimostrare risultati di espressività. Ad esempio \(FIN = \set{M \in \mathcal{M}}[\abs{M} < \infty]\) non è esprimibile in \(\FOL\) semplice, ovvero su linguaggio \(L = (\emptyset, \emptyset, \set{=})\). Supponiamo per assurdo che esista una formula \(\phi_{FIN}\) tale che \(\den{\phi_{FIN}} = FIN\). Consideriamo ora l'insieme parametrico di formule al variare di \(n\) definite come
\[\exists x_1, \ldots, x_n \bigwedge_{\substack{1 \leq i \leq n \\ 1 \leq j \leq n \\ i \neq j}} \neg(x_i = x_j)\]
Intuitivamente ognuna di queste formule esprime che esistono almeno \(n\) elementi distinti.

Consideriamo quindi \(\Gamma = \set{\phi_n}[n\in\N^*]\). Chiaramente questo insieme è numerabile. Sia ora \(\Sigma = \Gamma \cup \set{\phi_{FIN}}\). Ora se esiste un \(m\) che soddisfa \(\Gamma\) necessariamente deve essere infinito, poiché deve avere almeno \(n\) elementi distinti per ogni \(n\in \N\). Dunque \(\Sigma\) è insoddisfacibile perché qualsiasi modello di \(\Sigma\) dovrebbe essere sia finito che infinito. Ma ora considerando qualsiasi sott'insieme finito \(\Sigma_0\) di \(\Sigma\).

Chiaramente \(\Sigma_0\) è soddisfacibile, poiché basta un dominio finito di cardinalità maggiore o uguale al massimo indice \(n\) che compare in \(\Sigma_0\) per soddisfare tutte le formule di \(\Sigma_0\). Quindi abbiamo un insieme insoddisfacibile che però non ha nessun sottoinsieme finito insoddisfacibile, il che è assurdo per il teorema di compattezza.

Un altro esempio di particolare interesse per gli informatici è quello della non esistenza di una formula che caratterizza esattamente la classe dei grafi connessi. Chiaramente ciò non vuol dire che dato uno specifico grafo connesso non esista una formula che lo caratterizza o che esprima se quello specifico sia connesso, ma semplicemente che non esiste una formula che caratterizza tutti i grafi connessi e solo quelli.

Sia \(CONN = \set{M\in\mathcal{M}}[ M \text{ è un grafo connesso}]\). Supponiamo per assurdo che esista una formula \(\phi_{CONN}\) tale che \(\den{\phi_{CONN}} = CONN\). Per essere un grafo mi è sufficiente avere un linguaggio con un simbolo di relazione binaria \(E\) che rappresenta l'adiacenza tra i nodi con \(D\) che rappresenta i nodi. Per comodità aggiungiamo al linguaggio due costanti \(s\) e \(t\). Costruiamo quindi l'insieme di formule \(\Gamma = \set{\phi_n}[n\in\N \land n > 1]\) con
\[\phi_n \eqdef \neg\exists x_1, \ldots, x_n ( x_1 = s \land x_n = t \land \bigwedge_{1 \leq i < n} E(x_i, x_{i+1}))\]
Costruiamo quindi \(\Gamma' \eqdef \Gamma \cup \set{\forall x, \neg E(x,x)}\). Ogni formula \(\phi_n\) ci dice che non esiste un cammino di lunghezza \(n-1\) tra \(s\) e \(t\). Chiaramente preso \(M\) tale che \(M\models\Gamma'\) si ha che \(t^M\) non sarà raggiungibile da \(s^M\) e quindi \(M\) non sarà connesso. Dunque \(\Sigma \eqdef \Gamma' \cup \set{\phi_{CONN}}\) è insoddisfacibile. Ma ora se prendo un qualsiasi sottoinsieme finito \(\Sigma_0\) di \(\Sigma\) se questo non contiene \(\phi_{CONN}\) allora è banalmente soddisfacibile da qualsiasi grafo non connesso. Se invece \(\Sigma_0\) contiene \(\phi_{CONN}\) allora dalla finitezza esiste un indice massimo di formule di \(\Gamma\) che compare in \(\Sigma_0\), ovvero esiste \(n\) tale che \(\phi_n \in \Sigma_0\) e per ogni \(m > n\) si ha che \(\phi_m \not\in \Sigma_0\). A questo punto basta prendere un qualsiasi grafo in cui \(s\) e \(t\) sono connessi da un cammino di lunghezza strettamente maggiore di \(n\) per soddisfare tutte le formule di \(\Sigma_0\). Quindi anche in questo caso abbiamo un insieme insoddisfacibile che però non ha nessun sottoinsieme finito insoddisfacibile, il che è assurdo per il teorema di compattezza.

\subsection{Indecidibilità del prim'ordine}

Come anticipato precedentemente la logica del prim'ordine è indecidibile, e per mostrare ciò è necessario prima formalizzare precisamente tale idea. Riprendendo nozioni di teoria della complessità, nel dire che la logica del prim'ordine è indecidibile si intende che i problemi di decisione a essa associati, ovvero VALIDITY, SATISFIABILITY e LOGICAL CONSEQUENCE, sono indecidibili ovvero, riformulando tali problemi come linguaggi, e quindi la decisione come appartenenza a tali linguaggi, non esistono algoritmi che, in tempo finito, possano decidere l'appartenenza di un'arbitraria formula a tali linguaggi. Dunque, in altre parole, esistono formule per cui qualsiasi algoritmo che tenti di rispondere \qi{yes} o \qi{no} a seconda che la formula sia valida\footnote{Ovviamente riformulando il problema in un problema di appartenenza questo equivale a \q{la formula appartiene o meno al linguaggio delle formule valide}}, soddisfacibile o che sia conseguenza logica di un insieme di formule, non termina in tempo finito.

Come precedentemente accennato però è dimostrabile che questi problemi non sono \q{simmetrici} nel divergere. In particolare, esistono algoritmi che terminano in tempo finito per l'accettazione\footnote{Ovvero che rispondono \qi{yes} in tempo finito} di formule valide ma possono divergere per formule non valide, e algoritmi che terminano in tempo finito per il rifiuto\footnote{Ovvero che rispondono \qi{no} in tempo finito} di formule insoddisfacibili ma possono divergere per formule soddisfacibili. 

Per caratterizzare tali problemi è necessario introdurre alcune classi di complessità più fini, che non richiedono la totale decidibilità ma che si accontentino di uno dei due lati. Tali problemi sono detti semi-decidibili e appartengono alle classi RE (ricorsivamente enumerabili) e coRE (il complemento di RE).

Un insieme è ricorsivamente enumerabile se esiste una procedura che può enumerare tutti gli elementi dell'insieme. Invece coRE se il complemento dell'insieme può essere enumerato. 

Chiaramente avendo queste due procedure è semplice capire se il linguaggio è decidibile, poiché basta fare un passo alla volta alle enumerazioni di entrambe le procedure e vedere quale mi genera l'istanza.

In termini algoritmici, un linguaggio \(L\) è in RE se esiste un algoritmo che, dato un input \(x\), risponde \qi{yes} se \(x \in L\) e diverge altrimenti. Un linguaggio \(L\) è in coRE se esiste un algoritmo che, dato un input \(x\), risponde \qi{no} se \(x \not\in L\) e diverge altrimenti. Chiaramente, per essere decidibile, o ricorsivo, un linguaggio deve essere sia in RE che in coRE, ovvero deve esistere un algoritmo che risponde \qi{yes} se \(x \in L\) e \qi{no} se \(x \not\in L\) in tempo finito. Quello che andremo a mostrare in questa ultima sezione è che VALIDITY è in RE ma non è ricorsivo, e dunque per dualità e analogia si potrà concludere che SATISFIABILITY è in coRE ma non è ricorsivo, e che LOGICAL CONSEQUENCE è in RE ma non è ricorsivo.

Per mostrare ciò dovremo innanzitutto mostrare che VALIDITY è in RE, ovvero mostrare un algoritmo che almeno per le istanze positive del problema termini sempre con risposta \qi{yes}, e poi mostrare che ogni algoritmo che tenta di decidere completamente VALIDITY, e dunque anche di rispondere \qi{no} per le istanze negative porti a una contraddizione, ovvero che esistono formule per cui tale algoritmo diverge.

Iniziando con la prima parte, ovvero mostrare che VALIDITY è in RE, è sufficiente sfruttare i teoremi di completezza e correttezza. In particolare, vale \(\models \phi \iff \vdash \phi\). Riformulando quindi il problema in termini di appartenenza a linguaggi, si ha che \(\phi \in \Valid \iff \phi \in \set{\psi \in \FOL}[\vdash \psi]\). Dunque, per mostrare che VALIDITY è in RE, è sufficiente mostrare che l'insieme \(\set{\phi \in \FOL}[\vdash \phi]\) è in RE, ovvero che esiste una procedura per enumerarne tutti gli elementi. Per definizione di dimostrazione, una formula \(\phi\) è dimostrabile se esiste una derivazione di \(\phi\) a partire da \(\emptyset\). Essendo le derivazioni alberi finiti hanno tutti un altezza finita, e dunque è possibile iniziare a enumerare tutte le derivazioni di altezza \(0\), seguite da tutte le derivazioni di altezza \(1\), e così via. Tale enumerazione è ottenibile tramite una procedura poiché è sufficiente enumerare l'insieme finito di regole di inferenza e costruire tutte le possibili combinazioni di tali regole che portano a una derivazione di altezza \(n\) per ogni \(n\). Dunque, per ogni formula \(\phi\), se \(\phi\) è dimostrabile, allora esiste una derivazione di \(\phi\) a partire da \(\emptyset\) che sarà enumerata in un tempo finito, e dunque la procedura terminerà con risposta \qi{yes}.

Se si volesse mostrare adesso che VALIDITY è ricorsivo, si potrebbe pensare di mostrare che anche il complemento di VALIDITY, ovvero l'insieme delle formule non valide, è in RE, e dunque che esiste una procedura che termina con risposta \qi{yes} per ogni formula non valida. Un approccio intuitivo potrebbe essere quello di utilizzare lo stesso procedimento con le formule negate. Il problema sorge dal fatto che se \(\phi\) è valida, sicuramente \(\neg \phi\) è non valida, ma se \(\phi\) è non valida, non è detto che \(\neg \phi\) sia valida. In altre parole non tutte le formule non valide sono negazioni di formule valide, poiché è sufficiente che una formula sia soddisfacibile ma non valida per far si che la sua negata non sia valida. In altre parole l'insieme delle formule non valide è più grande dell'insieme delle negazioni di formule valide, e dunque non è sufficiente enumerare le negazioni di formule valide per enumerare tutte le formule non valide. 

A questo punto quindi si può provare a dimostrare che \(\overline{\text{VALIDITY}}\), ovvero il complemento di VALIDITY, non è in RE\@.

La teoria della complessità descrive diverse gerarchie di complessità per caratterizzare e confrontare problemi a seconda della loro difficoltà computazionale. Una nota gerarchia è la gerarchia polinomiale, che mette in relazione problemi di decisione risolvibili in tempo polinomiale con diversi modelli di computazione, come macchine di Turing deterministiche e non deterministiche. Una gerarchia che ci potrà essere utile è invece una delle tante che classifica i problemi indecidibili, ovvero la gerarchia aritmetica, in cui le proprietà sono classificate rispetto alle proprietà in teoria dei numeri. In particolare abbiamo \(\Sigma_1^0\) insieme delle formule per cui esiste una prova, \(\Pi_1^0\) insieme delle formule per cui non esiste una prova.

Un noto problema indecidibile è il problema della fermata di una macchina di Turing, ovvero data una qualsiasi macchina di Turing stabilire se essa termini in tempo finito o diverga. Come spesso accade in teoria della complessità, per mostrare che un problema è almeno difficile quanto un altro, e dunque in questo caso indecidibile, è sufficiente mostrare una riduzione, ovvero una procedura che, data un'istanza di un problema indecidibile, costruisca un'istanza del problema che vogliamo dimostrare essere indecidibile, in modo tale che la risposta alla nuova istanza sia \qi{yes} se e solo se la risposta alla vecchia istanza era \qi{yes}. Se ora il nuovo problema in esame fosse decidibile, allora sarebbe possibile decidere anche il problema indecidibile applicando a ogni sua istanza la riduzione seguita dalla procedura di decisione, portando a una contraddizione con l'indecidibilità nota del problema originale. 

Tali procedure di riduzione non possono essere arbitrarie, ma devono loro stesse appartenere a classi di complessità che siano più deboli della classe di complessità a cui il problema oggetto della dimostrazione dovrebbe appartenere. Se si vuole mostrare che un problema sia NP-completo, allora la procedura di riduzione non può essere a sua volta NP-completa, poiché altrimenti la complessità della procedura potrebbe risiedere nella procedura di riduzione stessa e non nel problema oggetto della dimostrazione. In questo caso, trattando di semi-decidibilità, è sufficiente che la procedura sia decidibile, ovvero che termini in tempo finito per ogni istanza, non potendo così \q{imputare} a essa un eventuale divergenza.

Andiamo dunque a mostrare tale riduzione, per poi dimostrare come essa verifichi la condizione di correttezza, ovvero che la formula costruita sia valida se e solo se la macchina di Turing termina. 

Iniziamo quindi a rivedere la definizione di macchina di Turing per poi mostrare come codificarla in una formula del prim'ordine.

\begin{definition}{Macchia di Turing}
  Una macchina di Turing è una tupla \(M = (Q, \Sigma, \delta, q_0, q_{n})\) dove:
  \begin{itemize}
    \item \(Q = \set{q_0, q_1, \ldots, q_t}\) è un insieme finito di stati;
    \item \(\Sigma = \set{\sigma_0, \sigma_1, \ldots, \sigma_m}\) insieme finito di simboli di nastro con \(\sigma_0 = \sqcup\) simbolo \q{blank};
    \item \(\delta: Q \times \Sigma \to Q \times \Sigma \times \{\leftarrow, \rightarrow\}\) funzione di transizione che, dato uno stato e un simbolo letto sul nastro, restituisce un nuovo stato, un simbolo da scrivere e una direzione in cui muovere la testina;
    \item \(q_0 \in Q\) è lo stato iniziale;
    \item \(q_{n} \in Q\) è lo stato di terminazione;
  \end{itemize}
\end{definition}

La semantica di questi modelli è data dalle \keyword{esecuzioni}, sequenze di \keyword{configurazioni} della macchina, ovvero istantanee dello stato complessivo della macchina in un dato momento. Formalmente:

\begin{definition}{Configurazione}
  Una configurazione \(C\) di una macchina di Turing \(M\) è una tripla \((k, i,t)\) con:
  \begin{itemize}
    \item \(k \in \set{0,\ldots,n}\) indice di stato, con \(q_k\in Q\);
    \item \(i \in \N\) indice di posizione della testina sul nastro;
    \item \(t: \N \to \Sigma\) funzione che associa a ogni cella del nastro il simbolo presente in quella cella.
  \end{itemize}

  L'insieme di tutte le possibili configurazioni di \(M\) è denotato con \(\mathcal{C}_M\).

  Indichiamo con \(C_0\) la configurazione iniziale di \(M\), ovvero la tripla \((0,0,t_0)\) con \(t_0 : \N \to \set{\sqcup}\).
\end{definition}

\begin{definition}{Esecuzione}
  Un esecuzione di una macchina di Turing \(M\) è una funzione \(\pi: \N \to \mathcal{C}_M\) che associa a ogni passo di esecuzione la configurazione della macchina \(M\) in quel passo.
  
  Si ha che \(\pi(0) = C_0\) e che, per ogni \(n\in\N\), se \(\pi(n) = (k, i, t_n)\) e \(\delta(q_k, t(i)) = (q_{k'}, \sigma_l, d)\) allora \(\pi(n +1) = (k', i', t_{n+1})\) con
  \begin{itemize}
    \item \(i' = i + 1\) se \(d = \rightarrow\) e \(i' = i - 1\) se \(d = \leftarrow\);
    \item \(t_{n+1}(j) = t_n(j)\) per ogni \(j \neq i\) e \(t_{n+1}(i) = \sigma_l\).
  \end{itemize}

  Un esecuzione è terminante se esiste \(j\in\N\) tale che \(\pi(j) = (n, i, t_j)\), ovvero se in qualche passo di esecuzione la macchina raggiunge lo stato di terminazione \(q_n\).
\end{definition}

Trattando macchine deterministiche, per ogni macchina di Turing \(M\) esiste una e una sola esecuzione \(\pi\) di \(M\). 
Quello che dobbiamo definire quindi è una formula \(\phi_M\) tale che\footnote{Data una tupla \(v\) e un indice \(1 \leq i \leq \abs{v}\), con \({(v)}_i\) indichiamo l'elemento di indice \(i\) della tupla \(v\). Ad esempio, se \(v = (k,i,t)\), allora \({(v)}_1 = k\), \({(v)}_2 = i\) e \({(v)}_3 = t\). Essendo nel nostro caso la terza componente una funzione, con \({(v)}_3(y)\) indichiamo l'applicazione della funzione alla variabile \(y\).}:
\[\models \phi_M \iff \exists i\st {(\pi_M(i))}_1 = n\]

Per lavorare in logica del prim'ordine, è necessario definire un linguaggio che ci permetta di rappresentare le configurazioni e le esecuzioni di \(M\). Intuitivamente, avendo descritto sia il nastro che l'esecuzione come funzioni con dominio \(\N\), possiamo costruire il linguaggio assumendo di avere tale dominio e definendo un interpretazione dei simboli coerente con la macchina di Turing. In particolare, sarà necessaria una costante \(0\) che rappresenterà l'inizio dei conteggi, sia di posizioni del nastro che di passi di esecuzione, due funzioni unarie \(s\) e \(p\) che intuitivamente verranno interpretate nelle funzioni di successore e predecessore, per permettere di esprimere lo spostamento della testina sul nastro e l'evoluzione dell'esecuzione nel tempo, e infine due insiemi di predicati binari, una per gli stati e un'altra per il nastro.

Per il resto della trattazione, per semplicità di notazione, fisseremo l'utilizzo della variabile \(x\) per descrivere i passi di esecuzione, e le variabili \(y\) e \(z\) per descrivere le posizioni sul nastro.

I predicati \(q_k\), con \(k \in \set{0,\ldots,n}\), che si occuperanno di descrivere le prime due componenti di una configurazione, saranno interpretati come segue:
\[\Interp(q_k) = \set{(x,y)\in\N^2}[{(\pi(x))}_1 = k \land {(\pi(x))}_2 = y ]\]
ovvero \(\Interp(q_k)\) è l'insieme di tutte le coppie \((x,y)\) tali che alla \(x\)-esima configurazione dell'esecuzione, la macchina si trovi nello stato \(q_k\) e la testina sia posizionata sulla cella \(y\).

I predicati \(\sigma_l\), con \(l \in \set{0,\ldots,m}\), che si occuperanno di descrivere il contenuto del nastro, saranno interpretati come segue:
\[\Interp(\sigma_l) = \set{(x,y)\in\N^2}[{(\pi(x))}_3(y) = \sigma_l] \]
ovvero \(\Interp(\sigma_l)\) è l'insieme di tutte le coppie \((x,y)\) tali che alla \(x\)-esima configurazione dell'esecuzione, il nastro contenga in posizione \(y\) il simbolo \(\sigma_l\).

A questo punto, avendo definito il linguaggio, non ci resta che definire un insieme di formule che descrivano il comportamento della macchina di Turing \(M\) in modo da poter mostrare la tesi.

Le prime formule che dobbiamo definire servono a imporre i vincoli necessari in modo che gli unici comportamenti descritti saranno coerenti con un esecuzione valida di \(M\). Iniziamo con l'assiomatizzare la proprietà di \(s\) e \(p\) di essere inverse.

\begin{equation}
  \phi_1 \eqdef \forall x, (p(s(x)) = x \land (\neg(x = 0) \to s(p(x)) = x))
\end{equation}

In aggiunta a questa formula sarebbe possibile aggiungere tutte le formule che assiomatizzano la teoria dei numeri naturali, ma non sarà necessario, poiché le formule che andremo a definire in seguito, che descrivono il comportamento di \(M\), implicheranno già tutte le proprietà necessarie per \(s\) e \(p\) per essere interpretati come funzioni di successore e predecessore su \(\N\).

Un altro vincolo che la macchina di Turing deve rispettare è che le celle di memoria sono atomiche, nel senso che, a ogni passo di esecuzione, ogni posizione del nastro può contenere solo e soltanto un simbolo. Per esprimere questo vincolo, è sufficiente la formula:

\begin{equation}
  \phi_2 \eqdef \forall x, \forall y, \left(\bigwedge_{i=0}^{m} \bigwedge_{\substack{j=0\\j \neq i}}^{m} \neg\left(\sigma_i\left(x,y\right) \land \sigma_j\left(x,y\right)\right)\right)
\end{equation}

Inoltre, a ogni passo di esecuzione, la macchina può trovarsi in un solo stato.

\begin{equation}
  \phi_3 \eqdef \forall x, \forall y, \forall z, \left(\bigwedge_{i=0}^{n} \bigwedge_{\substack{j=0\\j \neq i}}^{n} \neg\left(q_i\left(x,y\right) \land q_j\left(x,z\right)\right)\right)
\end{equation}

In questo caso è necessaria una doppia quantificazione per la posizione della testina sul nastro, poiché l'unico vincolo che si vuole imporre è l'unicità dello stato in cui si trova la macchina, e non l'unicità della posizione della testina.

Infine, per completare le proprietà di consistenza della macchina, ovvero quelle che escludono comportamenti anomali, dobbiamo garantire che, a ogni passo di esecuzione, anche essendo nello stesso stato, la macchina non possa avere la testina in due posizioni diverse, ovvero non possa leggere più celle contemporaneamente:

\begin{equation}
  \phi_4 \eqdef \forall x, \forall y, \forall z, \left(\neg(y = z)\to \bigwedge_{i=0}^{n} \neg\left(q_i\left(x,y\right) \land q_i\left(x,z\right)\right)\right)
\end{equation}

Per soddisfare queste formule è quindi necessario che, fissato un \(x\), esista un unico \(q_i\) che contiene coppie con \(x\) come primo elemento, e che tutte le coppie con \(x\) come primo elemento in tale \(q_i\) abbiano lo stesso secondo elemento, ovvero la posizione della testina. Inoltre, per ogni coppia \((x,y)\), questa può essere contenuta da un unico \(\sigma_l\), ovvero quello del simbolo presente in posizione \(y\) del nastro alla \(x\)-esimo passo di esecuzione.

In un certo senso, queste prima quattro formule impongono vincoli \q{statici} della macchina, nel senso che ne descrivono invarianti necessarie per mantenerne la coerenza, e di conseguenza escludere tutti i modelli che non rappresentano esecuzioni valide di \(M\). Le formule che andremo a definire adesso invece descrivono invece il comportamento dinamico di \(M\), ovvero come evolve l'esecuzione nel tempo, e dunque come evolvono le configurazioni della macchina da un passo di esecuzione al successivo. Intuitivamente, in queste formule sarà codificata la funzione di transizione \(\delta\). Iniziamo però dal codificare la configurazione iniziale. Come già detto la configurazione iniziale è \(C_0 = (0,0,t_0)\) con \(t_0 : \N \to \set{\sqcup}\). Tale configurazione dev'essere quella presente per prima nell'esecuzione, e quindi dev'essere quella associata a \(x = 0\). Per esprimere ciò, è sufficiente la formula:

\begin{equation}
  \phi_5 \eqdef q_0(0,0) \land \forall y, \sigma_0(0,y)
\end{equation}

Grazie ai vincoli di coerenza espressi prima, per assicurarci che all'inizio dell'esecuzione la macchina si trovi nella configurazione \(C_0\) è sufficiente descrivere come dev'essere fatta tale configurazione, poiché, grazie ai vincoli precedenti, sarà l'unica configurazione che soddisfa tali descrizioni.

A questo punto, prima di codificare propriamente la funzione di transizione, è necessario definire un ultimo vincolo di coerenza, che questa volta invece che essere locale assicura la coerenza tra due configurazioni consecutive. In particolare, coi vincoli precedenti abbiamo garantito che il nastro in ogni configurazione non abbia più di un simbolo nella stessa posizione e che la testina si trovi esattamente su una sola di esse, ma è necessario garantire anche che la posizione letta dalla testina sia l'unica che può essere modificata da un passo di esecuzione al successivo, e che tutte le altre posizioni del nastro rimangano inalterate. Per esprimere ciò, è sufficiente la formula:

\begin{equation}
  \phi_6 \eqdef \forall x, \forall y, \left(\bigwedge_{i=0}^{m} \left(\left(\sigma_i\left(x,y\right) \land \neg \sigma_i\left(s\left(x\right),y\right)\right) \to \left(\bigvee_{j=0}^{n} q_j\left(x,y\right)\right)\right)\right)
\end{equation}

Questa formula può sembrare complessa, ma in definitiva quello che richiede è che se in due istanti successivi (\(x\) e \(s(x)\)) il simbolo presente in una posizione \(y\) del nastro è diverso, allora la testina deve essere posizionata in \(y\) in uno qualsiasi degli stati. Questo vincolo, unito a quello che garantisce che in un istante di tempo la testina si trovi in una sola posizione, garantisce che l'unica posizione del nastro che può essere modificata da un passo di esecuzione al successivo è quella su cui è posizionata la testina. Se infatti da un passo al successivo più celle venissero modificate avremmo gli antecedenti di due di queste implicazioni in congiunzione verificati, ma al più uno dei conseguenti potrebbe essere verificato.

Infine, possiamo passare alla formula più complessa, che codifica la funzione di transizione \(\delta\). Per semplificarne la scrittura formeremo una formula aperta per ogni applicazione di \(\delta\), che poi andremo a congiungere e quantificare. Dunque presa un applicazione di \(\delta\) del tipo \(\delta(q_k, \sigma_l) = (q_{k'}, \sigma_{l'}, d)\), con \(d\) che può essere \(\rightarrow\) o \(\leftarrow\), definiamo le formule:
\begin{equation}
  \begin{aligned}
     \phi^{\delta^{\rightarrow}}_{k,l} (x,y) &\eqdef ((q_k(x,y) \land \sigma_l(x,y)) \to q_{k'}(s(x),s(y)) \land \sigma_{l'}(s(x),y))\\
     \phi^{\delta^{\leftarrow}}_{k,l} (x,y) &\eqdef ((q_k(x,y) \land \sigma_l(x,y)) \to q_{k'}(s(x),p(y)) \land \sigma_{l'}(s(x),y))
  \end{aligned}
\end{equation}
Dove \(k\) e \(l\) variano rispettivamente tra \(0\) e \(n\) e tra \(0\) e \(m\). Di queste chiaramente, per ogni coppia \((k,l)\) considereremo solo quella con direzione coerente a \(d\).

Possiamo dunque definire
\begin{equation}
  \phi^{\delta}_{k,l} (x,y) \eqdef \begin{cases}
    \phi^{\delta^{\rightarrow}}_{k,l} (x,y) & {(\delta(q_k, \sigma_l))}_3 = \rightarrow\\
    \phi^{\delta^{\leftarrow}}_{k,l} (x,y) & {(\delta(q_k, \sigma_l))}_3 = \leftarrow
  \end{cases}
\end{equation}

Per codificare quindi l'intera funzione di transizione sarà sufficiente utilizzare la formula:
\begin{equation}
  \phi_7 \eqdef \forall x, \forall y, \left(\bigwedge_{k=0}^{n} \bigwedge_{l=0}^{m} \phi^{\delta}_{k,l} (x,y)\right)
\end{equation}

Queste sette formule quindi descrivono precisamente la macchina di Turing \(M\) e il suo comportamento, nel senso che ogni modello di queste formule rappresenta un'esecuzione valida di \(M\), e ogni esecuzione valida di \(M\) è rappresentata da un modello di queste formule. A questo punto, per esprimere la proprietà di terminazione della macchina non ci serve altro che richiedere che se la congiunzione di queste condizioni è verificata esista un passo di esecuzione in cui lo stato di terminazione è raggiunto, ovvero che esista \(x\) e \(y\) tali che \(q_n(x,y)\) sia verificato. Dunque, la formula che codifica la terminazione macchina di Turing \(M\) è la seguente:
\begin{equation}
  \phi_M \eqdef \left(\bigwedge_{i=1}^{7} \phi_i\right) \to \exists x, \exists y, q_n(x,y)
\end{equation}

A questo punto, per completare la dimostrazione dell'indecidibilità di VALIDITY, è sufficiente mostrare che \(\phi_M\) è valida se e solo se \(M\) termina.

\begin{theorem}{Indecidibilità di VALIDITY}
  Sia \(M\) una macchina di Turing. Allora \(\models \phi_M\) se e solo se \(M\) termina, ovvero se esiste un \(i \in \N\) tale che \({(\pi_M(i))}_1 = n\).
\end{theorem}
\begin{proof}
  
  Per mostrare la doppia implicazione possiamo verificare separatamente il caso in cui \(M\) non termina e quello in cui termina.
  
  Iniziamo supponendo che \(M\) non termini, ovvero che non esista \(i \in \N\) tale che \({(\pi_M(i))}_1 = n\).
  Per mostrare che \(\not\models \phi_M\) è sufficiente mostrare che esiste un modello di \(\bigwedge_{i=1}^{7} \phi_i\) in cui \(\neg\exists x, \exists y, q_n(x,y)\) è verificato. Essendo \(M\) non terminante, esiste un'esecuzione \(\pi_M\) di \(M\) in cui non esiste \(i \in \N\) tale che \({(\pi_M(i))}_1 = n\). Dunque, prendendo come modello \(M = (\N,I)\) con \(I\) che interpreta i simboli del linguaggio come descritto precedentemente, si ha che \(M\) è un modello di \(\bigwedge_{i=1}^{7} \phi_i\) in cui \(\neg\exists x, \exists y, q_n(x,y)\) è verificato, e dunque \(M\) è un modello di \(\bigwedge_{i=1}^{7} \phi_i \land \neg\exists x, \exists y, q_n(x,y)\), e dunque \(\not\models \phi_M\).
  
  L'altro verso è invece più complesso. Si vuole mostrare che se \(M\) termina, ovvero se esiste \(i \in \N\) tale che \({(\pi_M(i))}_1 = n\), allora \(\models \phi_M\). Essendo che la proprietà di terminazione è finitaria per definizione, se \(M\) termina deve esistere un prefisso finito di \(\pi_M\) di lunghezza \(i\) tale che \({(\pi_M(i-1))}_1 = n\). Possiamo quindi analizzare per induzione questa sequenza finita, mostrando che la codifica della configurazione di ogni passo è conseguenza logica della codifica della macchina di Turing e che può inoltre essere espressa a partire dalla configurazione iniziale iterando la funzione di successore. Formalmente, se a passo \(k\) si ha \(\pi_M(k) = (j, i, t_k)\) allora vale:
  
  \begin{equation}
    \bigwedge_{h=1}^{7} \phi_h \models q_j(s^k(0), s^i(0)) 
  \end{equation}
  
  E per ogni posizione \(r\) del nastro al passo \(k\), con \(t_k(r) = \sigma_l\), vale anche
  \begin{equation}
    \bigwedge_{h=1}^{7} \phi_h \models \sigma_l(s^k(0), s^r(0))
  \end{equation}
  
  Chiaramente, dimostrato questo si ha la tesi, poiché valendo questa conseguenza logica per ogni passo di esecuzione, varrà in particolare per il primo passo in cui lo stato di terminazione è raggiunto. Di conseguenza si avrà che
  \begin{equation}
    \models \bigwedge_{h=1}^{7} \phi_h \models q_n(s^x(0), s^y(0))
  \end{equation}
  per qualche \(x\) e \(y\), e dunque \(\models \phi_M\) per introduzione di \(\exists\).
  
  Per dimostrare questa invariante verifichiamo innanzitutto il caso base, ovvero \(\pi_M(0) = C_0\). Questo caso è banalmente vero grazie a \(\phi_5\), interpretando chiaramente \(s^0\) come l'identità. Per il passo induttivo, supponiamo che valga per un passo \(k\), ovvero che \(\bigwedge_{h=1}^{7} \phi_h \models q_j(s^k(0), s^i(0))\) e che per ogni posizione \(r\) del nastro al passo \(k\), con \(t_k(r) = \sigma_l\), valga anche \(\bigwedge_{h=1}^{7} \phi_h \models \sigma_l(s^k(0), s^r(0))\). Dobbiamo mostrare che allora vale anche per il passo \(k+1\).
  
  Per definizione induttiva di \(\pi_M\) si ha che se \(\pi_M(k) = (j, i, t_k)\) e \(\delta(q_j, t_k(i)) = (q_{j'}, \sigma_{l'}, d)\) allora \(\pi_M(k +1) = (j', i', t_{k+1})\) con
  \begin{itemize}
    \item \(i' = i + 1\) se \(d = \rightarrow\) e \(i' = i - 1\) se \(d = \leftarrow\);
    \item \(t_{k+1}(r) = t_k(r)\) per ogni \(r \neq i\) e \(t_{k+1}(i) = \sigma_{l'}\).
  \end{itemize}
  
  Per ipotesi induttiva si ha che
  \[\bigwedge_{h=1}^{7} \phi_h \models q_j(s^k(0), s^i(0))\]
  e che per ogni posizione \(r\) del nastro al passo \(k\), con \(t_k(r) = \sigma_l\), vale anche
  \[\bigwedge_{h=1}^{7} \phi_h \models \sigma_l(s^k(0), s^r(0))\]
  
  Rispetto alla direzione, l'unica differenza tra le configurazioni d'arrivo è che in una la posizione della testina, codificata con l'esponente \(i\), è incrementata di \(1\) e nell'altra è decrementata di \(1\). Iniziamo assumendo che \(d = \rightarrow\).
  
  Grazie a \(\phi_7\) e alla definizione di \(\phi^{\delta}_{j,l}\), si ha che
  \[\bigwedge_{h=1}^{7} \phi_h \models (q_j(s^k(0),s^i(0)) \land \sigma_l(s^k(0),s^i(0))) \to q_{j'}((s(s^{k}(0)),s(s^{i})(0)) \land \sigma_{l'}(s(s^{k}(0)),s^{i}(0)))\]
  
  Chiaramente \(s(s^{k}(0)) = s^{k+1}(0)\) e \(s(s^{i}(0)) = s^{i+1}(0)\). Dunque essendo per ipotesi veri entrambi i congiunti di questa implicazione sarà anche vero
  \[\bigwedge_{h=1}^{7} \phi_h \models q_{j'}(s^{k+1}(0),s^{i+1}(0)) \land \sigma_{l'}(s^{k+1}(0),s^{i+1}(0))\]
  
  Inoltre, grazie a \(\phi_6\) si ha che tutte le posizioni diverse da \(i\) del nastro al passo \(k+1\) contengono lo stesso simbolo che avevano al passo \(k\), e dunque per ipotesi induttiva si ha che per ogni posizione \(r\) del nastro al passo \(k+1\), con \(t_{k+1}(r) = \sigma_l\), vale anche
  \[\bigwedge_{h=1}^{7} \phi_h \models \sigma_l(s^{k+1}(0), s^r(0))\]
  
  Venendo al caso di \(d = \leftarrow\), l'unica differenza si trova della definizione di \(\phi^{\delta}_{j,l}\), che in questo caso presenta la funzione \(p\). Possiamo notare però che grazie a \(\phi_1\) si ha che \(p(s^i(0)) = s^{i-1}(0)\) ottenendo esattamente la posizione della testina corretta al passo \(k+1\). 
\end{proof}